\documentclass[11pt]{article}

\usepackage[utf8]{inputenc}
\usepackage[T1]{fontenc}
\usepackage{geometry}
\usepackage{amsmath}
\usepackage{amssymb}
\usepackage{booktabs}
\usepackage{hyperref}
\usepackage{xurl}
\usepackage{xcolor}
\usepackage{listings}
\usepackage{enumitem}
\usepackage{parskip}
\usepackage{graphicx}
\graphicspath{{../images/}}

\geometry{margin=1in}

\lstset{
  language=Java,
  basicstyle=\ttfamily\small,
  keywordstyle=\color{blue},
  commentstyle=\color{gray},
  stringstyle=\color{teal},
  showstringspaces=false,
  columns=fullflexible,
  frame=single,
  framerule=0.2pt,
  breaklines=true
}

\setlist[itemize]{topsep=0.3em,itemsep=0.2em}
\setlist[enumerate]{topsep=0.3em,itemsep=0.2em}

% Simple per-section source line.
\newcommand{\notesources}[1]{\par\noindent{\small\color{gray}\textit{Sources:} \sloppy #1\par}}

% Unit-wise source strings for this subject (used by slides/notes).
% Path is relative to the lecture `latex/` folder (the compilation CWD).
% OOPS with Java (CSEG1044) — unit-wise sources
%
% These are short, student-facing reference strings intended to be used with:
%   - \slidesources{...}  (in beamer slides)
%   - \notesources{...}   (in lecture notes)
%
% Style inspiration: other subjects in My-Subjects/ use semicolon-separated sources
% and include an "accessed" date for web documentation.

\newcommand{\OOPSUnitISources}{Schildt, \textit{Java: A Beginner's Guide} (10e), McGraw-Hill (Java fundamentals + OOP basics).; Oracle Java Tutorials: Getting Started + Language Basics + Classes and Objects (accessed \today).; Java SE API docs (e.g., \texttt{String}, \texttt{Scanner}) (accessed \today).}

\newcommand{\OOPSUnitIISources}{Schildt, \textit{Java: The Complete Reference} (12e), McGraw-Hill (classes, inheritance, packages, interfaces).; Bloch, \textit{Effective Java} (3e), Addison-Wesley, 2018 (design choices; interfaces vs abstract classes).; Oracle Java Tutorials: Classes and Objects + Interfaces + Packages (accessed \today).; Java SE API docs (e.g., \texttt{Comparable}, \texttt{Runnable}) (accessed \today).}

\newcommand{\OOPSUnitIIISources}{Schildt, \textit{Java: The Complete Reference} (12e) (nested/inner classes, exceptions, threads, I/O).; Oracle Java Tutorials: Nested Classes + Exceptions + Concurrency + Basic I/O (accessed \today).; Java SE API docs (e.g., \texttt{Thread}, \texttt{AutoCloseable}) (accessed \today).}

\newcommand{\OOPSUnitIVSources}{Schildt, \textit{Java: The Complete Reference} (12e) (generics, lambdas, Swing, JDBC).; Bloch, \textit{Effective Java} (3e) (generics + lambdas).; Oracle Java Tutorials: Generics + Lambda Expressions + Swing + JDBC Basics (accessed \today).; Java SE API docs (e.g., \texttt{java.util.function}, \texttt{javax.swing}, \texttt{java.sql}) (accessed \today).}

\newcommand{\OOPSUnitVSources}{Schildt, \textit{Java: The Complete Reference} (12e) (Collections Framework + wrapper classes).; Oracle Java docs: Collections Framework overview (accessed \today).; Java SE API docs (\texttt{java.util} collections; wrapper classes in \texttt{java.lang}) (accessed \today).}

\newcommand{\OOPSUnitVISources}{Git SCM: \textit{Pro Git} (2e) (version control workflow), accessed \today.; GitHub Docs (repositories, issues, pull requests), accessed \today.; Oracle Java Tutorials: Swing + JDBC (accessed \today).; JUnit 5 User Guide (accessed \today).; Apache Maven Project: Getting Started (accessed \today).; Gradle User Manual: Getting Started (accessed \today).; UML class diagrams: UML 2.x overview/spec (accessed \today).}

\newcommand{\OOPSMiscWhyInterfacesSources}{Schildt, \textit{Java: The Complete Reference} (12e) (interfaces).; Bloch, \textit{Effective Java} (3e) (prefer interfaces where appropriate).; Oracle Java Tutorials: Interfaces (accessed \today).; Java SE API docs (e.g., \texttt{Comparable}, \texttt{Runnable}) (accessed \today).}



\title{Object Oriented Programming Lab (CSEG1044)\\Experiment-05\&06: Arrays + String Handling (Student Handout --- Before Submission)}
\author{Tofik Ali}
\date{\today}

\begin{document}
\maketitle
\tableofcontents

\section{About this handout}
This is a \textbf{pre-submission} handout.
It specifies what to implement and what to submit.

\section{Experiment 05 (Syllabus)}
\begin{itemize}[leftmargin=*]
  \item \textbf{Objective:} To work with one-dimensional and two-dimensional arrays.
  \item \textbf{Problem Statement:} Develop Java programs to perform operations on arrays.
  \item \textbf{Expected Outcome:} Students will be able to store and process data using arrays.
\end{itemize}

\section{Experiment 06 (Syllabus)}
\begin{itemize}[leftmargin=*]
  \item \textbf{Objective:} To perform string manipulation using Java String classes.
  \item \textbf{Problem Statement:} Implement string operations using built-in Java methods.
  \item \textbf{Expected Outcome:} Students will be able to manipulate textual data.
\end{itemize}

\section{What you must do (minimum requirements)}

\subsection{Task A --- Array operations (Exp-05)}
\begin{enumerate}[leftmargin=*]
  \item Create \texttt{MarksAnalysis.java}:
  \begin{itemize}[leftmargin=*]
    \item read \texttt{n} and then \texttt{n} marks into an \texttt{int[]},
    \item print min, max, sum, average,
    \item print \texttt{passCount} where marks $\ge 40$.
  \end{itemize}

  \item Create \texttt{MatrixAdd.java}:
  \begin{itemize}[leftmargin=*]
    \item read rows \texttt{r} and cols \texttt{c},
    \item read matrix \texttt{A} and matrix \texttt{B} of size \texttt{r x c},
    \item print \texttt{A + B} (element-wise sum).
  \end{itemize}
\end{enumerate}

\subsection{Task B --- String toolkit + palindrome (Exp-06)}
Create \texttt{StringToolkitDemo.java} and demonstrate:
\begin{enumerate}[leftmargin=*]
  \item \texttt{==} vs \texttt{equals()} using:
  \begin{itemize}[leftmargin=*]
    \item \texttt{new String("upes")}
    \item string literals like \texttt{"java"}
  \end{itemize}

  \item Immutability:
  \begin{itemize}[leftmargin=*]
    \item show that calling \texttt{toUpperCase()} does not change the original string unless stored.
  \end{itemize}

  \item Palindrome checker:
  \begin{itemize}[leftmargin=*]
    \item read a full line (sentence),
    \item ignore spaces and case,
    \item print whether it is a palindrome.
  \end{itemize}
\end{enumerate}

\section{Hints (read only if stuck)}
\begin{itemize}[leftmargin=*]
  \item Average should be computed using \texttt{sum / (double) n} (avoid integer division).
  \item For 2D arrays: \texttt{a.length} is number of rows; \texttt{a[i].length} is number of columns in row \texttt{i}.
  \item Use \texttt{equals()} for value comparison of strings; \texttt{==} checks references.
  \item Use \texttt{StringBuilder} when repeatedly appending in a loop.
\end{itemize}

\section{What to submit (MANDATORY)}
Submit a single zipped folder containing:
\begin{itemize}[leftmargin=*]
  \item \texttt{MarksAnalysis.java}
  \item \texttt{MatrixAdd.java}
  \item \texttt{StringToolkitDemo.java}
  \item sample output for all programs
  \item observations (5--10 lines): what you learned about arrays and string comparison
\end{itemize}

\notesources{\OOPSUnitISources}
\end{document}
